\expandafter\def\csname HistoricalNoteText15\endcsname{That it always has a completely definite course rests on its being part of the outer boundary of a region in which $T$ has a definite sign. I have chosen the positive sign, which in itself is entirely arbitrary. Understood in this way, for example, the line entering at $1$ continues through $11$ to $10$. As part of the boundary of a region in which $T$ is negative, the line $1.11$ would have to be continued to $2$. If, on the other hand, one speaks only of a line on which $T=0$, without regarding it as part of the boundary of a particular region, $11.9$ could more readily be considered the natural continuation of $1.11$. The point of view chosen here distinguishes my present procedure from that of $1799$, and contributes substantially to simplifying the proof.}
